\documentclass[12pt]{article} % Увеличенный базовый шрифт
\usepackage[utf8]{inputenc}
\usepackage[T2A]{fontenc}
\usepackage[english,russian]{babel}
\usepackage{amsmath}
\usepackage[left=1.5cm, right=1.5cm, top=1.5cm, bottom=1.5cm]{geometry}

\title{\begin{center}Вопрос по выбору:\end{center}
Уравнения моделей неидеальных газов}
\author{Струков О. И. \\ Б04-404}
\date{}

\begin{document}

    \pagenumbering{gobble}
    \maketitle
    \newpage
    \pagenumbering{arabic}


\section*{Другой метод введения поправки на силы притяжения между молекулами. Уравнение Дитеричи}

\subsection*{1. }

Влияние пристеночного молекулярного слоя на уравнение состояния можно учесть другим способом, который приведется ниже. Отвлечёмся сначала от молекулярных сил притяжения, с которыми стенка может действовать на молекулы газа. Будем предполагать, что стенка действует на молекулы газа только при столкновениях с ними. Как выяснено в предыдущем параграфе, молекулы пристеночного слоя подвергаются действию результирующей силы, направленной внутрь газа. Вследствие этого концентрация молекул в пристеночном слое должна убывать при приближении к стенке в соответствии с формулой Больцмана:

\begin{equation}
n = n_\infty e^{-\frac{U}{kT}}, \tag{99.1}
\end{equation}

где $U$ — потенциальная энергия молекулы. Этим и объясняется уменьшение давления газа на стенку. Энергия $U$ является функцией расстояния $x$ молекулы от стенки. Она максимальна у самой стенки и быстро убывает с возрастанием $x$. Ее значение на бесконечности $U_\infty$ условимся считать равным нулю. (Бесконечными считаются расстояния больше радиуса сферы молекулярного действия.) Давление на стенку определяется пристеночной концентрацией $n_\infty$, а не внутренней концентрацией $n$. Считая молекулы точечными, можем написать, как для идеального газа:

\begin{equation}
P = n_\infty kT = n_\infty kTe^{-\frac{U_0}{kT}}, \tag{99.2}
\end{equation}

где $U_0$ — значение потенциальной энергии молекулы у стенки сосуда. Если выполнено условие $U_\infty \ll kT$, то $e^{-\frac{U_\infty}{kT}} \approx 1 + \frac{U_\infty}{kT}$, а потому

\begin{equation}
P + n_\infty U_\infty = n_\infty kT. \tag{99.3}
\end{equation}

Сила притяжения, действующая на пристеночную молекулу, а с ней и потенциальная энергия $U_\infty$, пропорциональны концентрации молекул газа $n_\infty = N/V$. Поэтому можно написать $U_\infty = \alpha n_\infty$, где $\alpha$ — постоянная для рассматриваемого газа величина. Предыдущая формула преобразуется к виду $(P + \alpha n_\infty^2) = n_\infty kT$, или

\begin{equation}
\left(P + \frac{\alpha}{V}\right)V = RT, \tag{99.4}
\end{equation}

если ввести новую постоянную $a = \alpha N^2$. Этот результат совпадает с (98.5). Заметим еще, что через постоянную $a$ потенциальная энергия $U_\infty$ выражается формулой

\begin{equation}
U_\infty = \frac{a}{NV}. \tag{99.5}
\end{equation}

\subsection*{2. Молекулы пристеночного слоя притягиваются не только молекулами газа, но и молекулами стенки}

Это должно привести к уплотнению пристеночного слоя и к увеличению давления газа на стенку. Получается парадоксальный результат: давление газа на стенку и постоянная $a$ в уравнении Ван-дер-Ваальса должны зависеть от материала стенки. Этот результат не согласуется с опытом, а вывод его неверен. Разберемся в этом вопросе.

Сила $F$, действующая на единицу площади стенки из-за ударов нелетящих молекул, действительно возрастает. Но это не единственная сила, действующая на стенку. Стенка притягивает газ. По третьему закону Ньютона газ притягивает стенку с равной, но противоположно направленной силой. Обозначим величину ее $F_1$. Чтобы получить давление газа на стенку, силу $F_1$ надо вычесть из $F$. Покажем, что в пределах точности расчета, принятой при выводе уравнения Ван-дер-Ваальса, результат не будет зависеть от материала стенки.

Пусть $U'(x)$ — потенциальная энергия молекулы в поле сил притяжения стенки. Формулу Больцмана теперь надо писать в виде

$$
n = n_\infty e^{-\frac{U + U'}{kT}}.
$$

Поэтому вместо (99.3) мы придем к соотношению

$$
F = n_\infty kT - n_\infty (U_0 + U'_0).
$$

Вычислим теперь силу $F_1$. Стенка притягивает молекулу с силой $f_1$, абсолютная величина которой равна $f_1 = dU'/dx$. Число молекул в слое единичной площади и толщины $dx$ есть $ndx$. Следовательно,

$$
F_1 = \int\limits_0^\infty n f_1 dx = \int\limits_0^\infty n \frac{dU'}{dx} dx = \int\limits_{U'_0}^0 ndU'.
$$

В принятом приближении зависимость концентрации $n$ от координаты $x$ учитывать не надо, а потому $F_1 = -n_\infty U'_0$. Вычитая эту величину из $F$, найдем

$$
P = F - F_1 = n_\infty kT - n_\infty U_0,
$$

что совпадает с (99.3). Таким образом, давление $P$ не зависит от материала стенки.

Обобщить это доказательство на случай произвольно больших плотностей газа, положив в основу рассмотрение механизма явления, довольно затруднительно. Однако в этом и нет необходимости. Независимость давления газа на стенку от материала стенки можно доказать на основе общих соображений. Рассмотрим закрытый цилиндрический сосуд, наполненный газом, противоположные основания которого $AB$ и $CD$ сделаны из различных материалов. Допустим, что установилось состояние равновесия. Если бы давления на стенки $AB$ и $CD$ были разными, то сосуд пришел бы в движение, и равновесие, вопреки предположению, было бы невозможно.



3. Вернемся к формуле (99.2), но не будем аппроксимировать показательную функцию линейно. Тогда с учетом соотношения (99.5) получим

$$
P = n_\infty kTe^{-\frac{a}{RTV}}.
$$

Введем сюда поправку на конечный объем молекул совершенно так же, как это делалось в предыдущем параграфе. Тогда

\begin{equation}
P(V - b) = RT e^{-\frac{a}{RTV}}. \tag{99.6}
\end{equation}

Это уравнение Дитеричи. В пределе, когда $b \ll V$, а $a \ll RTV$, уравнение (99.6) переходит в уравнение Ван-дер-Ваальса. Чтобы выполнить такой предельный переход, аппроксимируем в формуле (99.6) показательную функцию линейной. Получим

$$
P = \frac{RT}{V - b} - \frac{a}{V(V - b)}.
$$

В последнем слагаемом величиной $b$ следует пренебречь. Это приводит к ошибке второго порядка по малым поправкам $a$ и $b$. Величинами же такого порядка при выводе уравнения Ван-дер-Ваальса мы пренебрегали. Таким образом, мы снова возвращаемся к уравнению (98.10).

Уравнение Дитеричи является таким же полуэмпирическим уравнением, что и уравнение Ван-дер-Ваальса. Оба уравнения можно считать теоретически обоснованными только при выполнении условий, довольно затруднительно. Однако уравнение Ван-дер-Ваальса благодаря своей простоте и ясному физическому смыслу входящих в него постоянных $a$ и $b$ остается наиболее распространенным уравнением для анализа качественного поведения реальных газов и жидкостей. Приведем некоторые из наиболее известных уравнений состояния:

Уравнение Бертло:
\begin{equation}
\left(P + \frac{a}{TV^2}\right)(V - b) = RT. \tag{99.7}
\end{equation}

Уравнение Клаузиуса:
\begin{equation}
\left(P + \frac{a}{T(V + c)^2}\right)(V - b) = RT. \tag{99.8}
\end{equation}

Уравнение Камерлинга — Оннеса:
\begin{equation}
PV = RT\left(1 + \frac{B_2}{V} + \frac{B_3}{V^2} + \ldots\right), \tag{99.9}
\end{equation}

где $B_2$, $B_3$, ... называются вторым, третьим и последующими вириальными коэффициентами. Они являются функциями температуры. Таким уравнением мы уже пользовались в § 33 при приведении шкалы газового термометра к термодинамической шкале. Там же было показано, что уравнение состояния всякого газа может быть приведено к виду (99.9). Однако уравнение (99.9) получает конкретное содержание только после того, как будут найдены выражения для входящих в него виряльных коэффициентов, как функций температуры.

Уравнения Бертло и Клаузиуса отличаются от уравнения Ван-дер-Ваальса поправками, вводимыми чисто эмпирически. Уравнение Бертло вблизи критической точки не имеет никаких преимуществ по сравнению с уравнением Ван-дер-Ваальса. Зато при умеренных давлениях оно лучше согласуется с опытом. Уравнение Клаузиуса точнее уравнения Ван-дер-Ваальса, поскольку оно содержит третью эмпирическую постоянную $c$. Благодаря этому с помощью уравнения Клаузиуса можно учесть отклонения от закона соответствующих состояний.


\end{document}